\section[Forced interface lemmas]{Forced interface lemmas under the \texorpdfstring{tick\,+\,CAP}{tick+CAP} spine}
\label{app:forced_interface_lemmas}

This appendix records several short interface lemmas in a form suitable for audit.
It introduces \emph{no new axioms} beyond the two declared primitives of the paper: tick as the executed input stream (Axiom~\ref{ax:readout_sequentiality}) and CAP as the unique closure/selection rule on explicit finite candidate families (Axiom~\ref{ax:cap}).
The purpose is to make explicit the sense in which several ``interface conventions'' used in the main text are \emph{forced} (or uniquely selected as minimal nontrivial choices) once one commits to the tick\,+\,CAP discipline.

\subsection{Minimal coarse locking: one bit per independent parameter}
\label{subsec:forced_minimal_coarse_lock}

\begin{lemma}[Minimal nontrivial coarse binning forces one bit per independent parameter]
\label{lem:minimal_one_bit_per_parameter}
Let a protocol claim to \emph{coarsely distinguish} $k$ independent parameters in a single observation, in the minimal nontrivial sense that for each parameter there exist at least two disjoint bins whose membership is distinguishable in the readout.
Then any single-shot readout alphabet that supports such a coarse distinction must have size at least $2^k$.
In particular, if the single-shot readout alphabet is binary length-$m$ words $\Omega_m=\{0,1\}^m$, then necessarily $m\ge k$.
\end{lemma}
\begin{proof}
For each parameter choose two distinguishable bins and label them by $\{0,1\}$.
Independence means that all $2^k$ bin-combinations are admissible joint coarse states.
Therefore the readout alphabet must have at least $2^k$ distinct outcomes.
If the alphabet is $\Omega_m$, then $|\Omega_m|=2^m\ge 2^k$, so $m\ge k$.
\end{proof}

\begin{corollary}[Single-window coarse rigid-frame budget]
\label{cor:forced_rigid_frame_budget}
In a rigid-frame display dictionary in bulk dimension $d$, a pose is an element of the Euclidean group $SE(d)$ with
\[
\dim SE(d)=\frac{d(d+1)}{2}.
\]
Under the minimal nontrivial coarse-binning convention of Lemma~\ref{lem:minimal_one_bit_per_parameter}, a single length-$m$ window can support a single-shot coarse rigid-frame display only if
\[
m\ge \dim SE(d)=\frac{d(d+1)}{2}.
\]
\end{corollary}
\begin{proof}
Apply Lemma~\ref{lem:minimal_one_bit_per_parameter} with $k=\dim SE(d)$.
The dimension formula for $SE(d)$ is standard; see, e.g., \cite{HallLieGroups2015}.
\end{proof}

\paragraph{Relation to the main text.}
Section~\ref{subsec:6dof_lock} uses exactly this minimal nontrivial convention at the anchor $m=6$ and then applies CAP to select the maximal admissible bulk dimension, yielding $d=3$ (Proposition~\ref{prop:bulk_dimension_from_anchor}).
The quantization/metric-entropy scaling used there for finer accuracy is standard \cite{GrayNeuhoff1998Quantization,CoverThomas2006}.

\subsection{Unitarity forces compact internal redundancy groups}
\label{subsec:forced_unitary_compact}

\begin{proposition}[Probability-preserving internal redundancy is compact (finite-dimensional case)]
\label{prop:unitary_implies_compact_redundancy}
Assume a continuum modeling dictionary in which local internal redundancy acts on a finite-dimensional complex Hilbert space by transformations that preserve transition probabilities between rays.
Then the connected component of the internal redundancy group is (projectively) unitary and is therefore compact (up to finite quotients).
\end{proposition}
\begin{proof}
By Wigner's theorem, any bijection of rays preserving transition probabilities is implemented by a unitary or antiunitary operator on the underlying Hilbert space, uniquely up to phase \cite{Wigner1931,StreaterWightman1989,WeinbergQFT1995}.
Antiunitary operators form a disconnected component, so any connected internal redundancy group is represented (projectively) by unitary operators.
Thus the redundancy group embeds as a (closed) subgroup of the compact projective unitary group $PU(N)=U(N)/U(1)$, hence is compact.
\end{proof}

\paragraph{Relation to the main text.}
In the interface language of this paper, ``gauge redundancy'' is the freedom to relabel local readout bases without changing observable overlap/probability data.
Under a standard continuum dictionary that represents this freedom by internal unitary rotations, compactness is therefore not an extra primitive but a consequence of probability preservation in finite-dimensional local descriptions.
This is the compactness input used in Proposition~\ref{prop:channel_to_gauge}.

\subsection{Three commuting channels force a three-factor redundancy structure}
\label{subsec:forced_three_factor}

\begin{lemma}[Independent local redundancy across commuting channels factorizes]
\label{lem:three_channel_factorization}
Assume protocol mismatch certificates decompose into three commuting channel components and that, in the interface dictionary, each channel admits an independent local basis redundancy.
Then the minimal local redundancy group factorizes as a direct product of three channelwise redundancies.
\end{lemma}
\begin{proof}
Since channels commute, a mismatch certificate can be compared channel-by-channel.
If the basis redundancy is independent in each channel, local gauge changes act independently on each component.
Therefore the channelwise covariance law is a direct-product action, and the minimal redundancy group is the direct product of the three channel groups.
\end{proof}

\paragraph{Relation to the main text.}
This is the structural content of Proposition~\ref{prop:three_factor_compensation} in Section~\ref{subsec:gauge_as_compensation}.
Once a three-factor structure is fixed, CAP-minimal selection within the compact-factor candidate family yields the Standard Model triple up to finite quotients under the declared factor complexity label (Proposition~\ref{prop:channel_to_gauge} and Lemma~\ref{lem:su2_su3_unique_by_dim}); a bounded sensitivity sweep across alternative labels is recorded in Appendix~\ref{app:gauge_complexity_sensitivity}.

\subsection{Why the minimal chiral closure selects a sterile \texorpdfstring{$\nu_R$}{nuR}}
\label{subsec:forced_nur}

\begin{proposition}[Minimal anomaly-neutral one-multiplet-per-generation closure forces a sterile singlet]
\label{prop:nur_forced_by_minimal_anomaly}
Work under the PDG convention $Q=T_3+Y$.
Consider extending the Standard Model chiral fermion multiplet content by adding \emph{exactly one} additional chiral multiplet per generation, with the requirement that:
(i) local gauge and mixed gravitational anomalies remain canceled, and
(ii) the global $SU(2)$ consistency condition is preserved.
Then the added multiplet must be a gauge singlet with hypercharge $Y=0$.
Equivalently (up to charge conjugation), the unique minimal closure is a sterile right-handed neutrino $\nu_R$ per generation.
\end{proposition}
\begin{proof}
The Standard Model anomaly-cancellation identities under $Q=T_3+Y$ are standard \cite{BouchiatIliopoulosMeyer1972,PeskinSchroeder1995,PDGRPP2024}.
If the added multiplet carries nontrivial $SU(3)$ charge, it contributes to $SU(3)^3$ and $SU(3)^2U(1)$ anomaly sums; with only one added multiplet per generation there is no compensator, so anomaly cancellation forces it to be an $SU(3)$ singlet.

If the added multiplet is an $SU(2)$ doublet (or any half-integer isospin representation), the parity of the number of $SU(2)$ doublets changes and can violate the global $SU(2)$ anomaly constraint \cite{Witten1982SU2Anomaly}.
With only one additional multiplet per generation and no compensator, preserving the global consistency condition forces the added multiplet to be an $SU(2)$ singlet.

Thus the only remaining possible charge is hypercharge.
The mixed gravitational--hypercharge anomaly is proportional to $\sum Y$ and the cubic $U(1)_Y^3$ anomaly to $\sum Y^3$ over left-handed Weyl fields with multiplicities; the Standard Model sums vanish per generation \cite{BouchiatIliopoulosMeyer1972,PeskinSchroeder1995}.
Adding a single $SU(3)\times SU(2)$ singlet multiplet contributes $Y$ and $Y^3$.
To preserve both anomaly cancellations without additional compensating matter, one must have $Y=0$.
\end{proof}

\paragraph{Relation to the main text.}
Section~\ref{sec:sm_labeling_closure} uses $\nu_R$ as the minimal anomaly-neutral closure of the $18$ cyclic field-level targets and treats it as an interface choice audited by standard consistency constraints (Proposition~\ref{prop:anomaly_nur}).
Proposition~\ref{prop:nur_forced_by_minimal_anomaly} records the stronger uniqueness statement within the explicit minimal candidate family ``one extra multiplet per generation''.

\subsection{Scalar absence at the minimal alphabet is forced by the closed \texorpdfstring{$21$}{21}-type contract}
\label{subsec:forced_scalar_absence}

\begin{proposition}[No additional primitive stable type remains for a Higgs label at \texorpdfstring{$m=6$}{m=6}]
\label{prop:higgs_not_in_x6_forced}
At the anchor $(m,n)=(6,3)$, the stable sector splits as $X_6=X_6^{\mathrm{cyc}}\sqcup X_6^{\mathrm{bdry}}$ with $|X_6^{\mathrm{cyc}}|=18$ and $|X_6^{\mathrm{bdry}}|=3$.
Under the closed interface contract adopted in this paper---cyclic types label the $18$ field-level chiral multiplets and boundary types label the three gauge-factor connection classes---there is no remaining element of $X_6$ that can be assigned as an additional primitive stable-type label for a Higgs-like scalar.
Accordingly, any scalar sector must enter either as (i) an EFT-level completion field in the continuum dictionary, or (ii) a protocol-emergent observable supported by resolution uplift and coarse graining.
\end{proposition}
\begin{proof}
The split cardinalities are theorem-level facts at $m=6$ (Section~\ref{sec:folding_core}).
The closed labeling map assigns the $18$ cyclic types to the $18$ chiral multiplets and assigns the $3$ boundary types to the three gauge factors (Theorem~\ref{thm:labeling_unique}).
Therefore all $21$ stable types are already consumed by the closed chiral/gauge interface at the anchor, leaving no unused stable label in $X_6$.
\end{proof}

\paragraph{Relation to the main text.}
This is the audit-level content behind the Higgs-status remarks in Section~\ref{sec:sm_labeling_closure} (Remark~\ref{rem:higgs_not_in_21}) and the scalar-sector closure statements in Section~\ref{subsec:scalar_sectors_interface}.


